%!TEX root = ../DA_MainDocument.tex
\chapter{Systemarchitektur \& technische Übersicht}\label{chapter:architektur}

Dieses Kapitel beschreibt die Architektur von Skyline und die technische Umsetzung
jener Features, die in den Implementierungskapiteln~5--8 nicht behandelt werden:
Flugverwaltung, Dokumente, Notizen, Checklisten, Statistiken, Company-Features und
Gamification. Die vier Forschungsmodule -- Import (Kapitel~\ref{chapter:import}),
Datenverwaltung (Kapitel~\ref{chapter:datenverwaltung}), Weltkarte
(Kapitel~\ref{chapter:weltkarte}) und Benachrichtigungen
(Kapitel~\ref{chapter:notifications}) -- werden hier nur referenziert.

Skyline wurde mit React Native und Expo als mobile Cross-Platform-App umgesetzt, im Projektzeitraum primär für iOS entwickelt und getestet und ist im \textbf{Apple App Store} verfügbar. Eine Android-Variante existiert, wurde jedoch weniger umfangreich validiert.

% ============================================================
\section{Systemarchitektur}
% ============================================================

Die Architektur folgt einem Client-Server-Modell: Die mobile App kommuniziert
ausschließlich über HTTPS/REST mit Supabase als Backend. Eine zentrale
\textbf{Service-Schicht} entkoppelt UI-Komponenten von der Datenhaltung und kapselt
alle Supabase-Zugriffe, Dokumentenoperationen, Reminder-Planung und
Import-Prozesse. Sämtliche Sicherheitslogik liegt in der Datenbank selbst (RLS);
die Service-Schicht ist bewusst kein Sicherheitsmechanismus, sondern ein
Wartbarkeits-Layer.

\begin{figure}[H]
  \centering
  \vspace{1.5em}
  \begin{tikzpicture}[
    archbox/.style={draw, rounded corners, align=center, minimum width=3cm, minimum height=1cm, font=\small},
    arrow/.style={-Latex, thick}
  ]
    \node[archbox] (client) at (0, 3.8) {Mobile Client\\React Native + Expo\\(primär iOS getestet)};
    \node[archbox] (services) at (0, 1.2) {Service-Schicht\\services/supabase.ts\\documentService, ocr, company};
    \node[archbox] (supabase) at (6.2, 2.5) {Supabase\\Auth + Postgres\\RLS, Storage, Realtime};
    \node[font=\scriptsize] (store) at (-1.8, 2.5) {store/index.ts};
    \draw[arrow] (client.south) -- node[right, font=\scriptsize] {API-Calls} (services.north);
    \draw[arrow] (services.east) -- node[above=3mm, pos=0.5, font=\scriptsize, fill=white, inner sep=2pt] {HTTPS/REST} (supabase.west);
    \draw[->, dashed] (client.west) -- ++(-0.3,0) node[left, font=\scriptsize]{UI-Screens};
  \end{tikzpicture}
  \vspace{2.5em}
  \caption{Systemarchitektur von Skyline: Client, Service-Schicht und Supabase-Backend}
  \label{fig:architektur}
\end{figure}

\subsection{Client}

Der Client basiert auf \textbf{React Native} und \textbf{Expo} \cite{reactNative,expoDocs}.
\textbf{Expo Router} übernimmt die dateibasierte Navigation; der globale Anwendungszustand wird mit
\textbf{Zustand} in einem zentralen Store (\texttt{store/index.ts}) verwaltet -- Flugdaten,
Nutzerprofil und Einstellungen liegen dort zentral vor. Native iOS-APIs für Kamera,
Fotoauswahl und Datei-Picker werden über die jeweiligen Expo-Module eingebunden. Der
iOS-Build wird über \textbf{EAS Build} erzeugt und im \textbf{Apple App Store} bereitgestellt \cite{easBuild}.

\subsection{Backend}

\textbf{Supabase} liefert Authentifizierung (Supabase Auth), eine Postgres-Datenbank,
optionale Realtime-Subscriptions und Row-Level Security \cite{supabasePlatform,supabaseAuth,supabaseRLS}.
Alle operativen Tabellen -- darunter \texttt{user\_flights}, \texttt{user\_notes},
\texttt{user\_checklists}, \texttt{flight\_documents}, \texttt{companies},
\texttt{company\_members}, \texttt{notifications} und \texttt{user\_achievements} --
sind durch RLS-Policies geschützt. Details zum Datenbankschema
und zur RLS-Architektur sind in Kapitel~\ref{chapter:datenverwaltung} dokumentiert.

\subsection{Storage}

Dokumente und Bilder werden im privaten Storage-Bucket \texttt{flight-documents} abgelegt
und ausschließlich über zeitlich begrenzte Signed URLs bereitgestellt \cite{supabaseStorage}.
Kein öffentlicher Direktzugriff ist möglich. Metadaten (Dateiname, MIME-Type, Größe,
Storage-Pfad) werden in der Tabelle \texttt{flight\_documents} persistiert. Details zu
Upload-Ablauf, URL-Refresh und Rollback-Logik sind in Kapitel~\ref{chapter:datenverwaltung} beschrieben.

% ============================================================
\section{Technologien \& externe APIs}
% ============================================================

\subsection{Frontend-Stack}

\begin{table}[H]
  \centering
  \caption{Frontend-Stack von Skyline}
  \label{tab:frontend_stack}
  \begin{tabular}{p{4.2cm}p{8.5cm}}
  \toprule
  \textbf{Technologie} & \textbf{Zweck} \\
  \midrule
  React Native + Expo & Cross-Platform-Basis, primär für iOS getestet \cite{reactNative,expoDocs} \\
  Expo Router & Dateibasierte Navigation \\
  Zustand & Globaler State-Management-Store \\
  TypeScript & Einheitliche Typsicherheit im gesamten Codebase \\
  React Native Maps & Kartenvisualisierung (MapKit auf iOS) \cite{reactNativeMaps} \\
  EAS Build & Build- und Deployment-Pipeline (App Store) \cite{easBuild} \\
  \bottomrule
  \end{tabular}
\end{table}

\subsection{Backend-Stack}

\begin{table}[H]
  \centering
  \caption{Backend-Stack von Skyline}
  \label{tab:backend_stack}
  \begin{tabular}{p{4.2cm}p{8.5cm}}
  \toprule
  \textbf{Technologie} & \textbf{Zweck} \\
  \midrule
  Supabase Auth & Authentifizierung (JWT, Session-Management) \\
  Postgres + RLS & Relationale Datenhaltung mit zeilenbasierter Zugriffssteuerung \cite{supabaseRLS} \\
  Supabase Storage & Datei-Storage mit Signed URLs \cite{supabaseStorage} \\
  Supabase Realtime & Optionale Live-Subscriptions \\
  \bottomrule
  \end{tabular}
\end{table}

\subsection{Externe APIs}

\textbf{Aviationstack} liefert Airport-Stammdaten (IATA-Codes, ICAO-Codes, Koordinaten,
Stadt, Land), die bei Import und manueller Eingabe zur Validierung und
Distanzberechnung herangezogen werden \cite{aviationstack}.

\textbf{OCR.space} ist der eingesetzte OCR-Dienst für die Texterkennung aus Bildern und
PDFs \cite{ocrSpace}. Anfragen werden als HTTP-POST mit Base64-kodiertem Bildinhalt an
\url{https://api.ocr.space/parse/image} geschickt; als Engine wird OCREngine~2 (präziser
für gedruckten Text) verwendet. Der API-Key wird über die Umgebungsvariable
\texttt{EXPO\_PUBLIC\_OCR\_SPACE\_API\_KEY} konfiguriert. Details zur Parsing-Pipeline sind in
Kapitel~\ref{chapter:import} dokumentiert.

\textbf{Apple MapKit} (via React Native Maps) stellt die Kartendarstellung auf iOS bereit.
Auf Android wird die Google Maps API als sekundäre Plattform genutzt. Die Berechnung
geodätischer Routen (Great-Circle-Polylines) erfolgt appintern; externe
Routing-APIs werden nicht verwendet. Details zur Kartenimplementierung sind in
Kapitel~\ref{chapter:weltkarte} dokumentiert.

\textbf{E-Mail-Import:} Die Architektur sieht OAuth-verknüpfte E-Mail-Accounts
(Gmail, Outlook, Yahoo, IMAP) über die Tabellen \texttt{email\_accounts} und
\texttt{email\_imports} in Supabase vor; ein dedizierter E-Mail-Parser-Service ist
nicht implementiert. Umgesetzt ist der OCR-basierte Import von Buchungsbestätigungen
aus Screenshots und PDFs (vom Nutzer hochgeladen); E-Mail-Import ist als Erweiterung
vorgesehen.

% ============================================================
\section{Flugverwaltung}
% ============================================================

Flüge sind die Kernentität von Skyline. Alle anderen Module -- Notizen, Checklisten,
Dokumente, Reminder -- sind über \texttt{flight\_id} an einen Flug gebunden.

\subsection{Datenfluss: UI $\rightarrow$ Store $\rightarrow$ Service}

Es gibt keinen separaten Middleware-Layer. Die UI-Screens (\texttt{add-flight-manual.tsx},
\texttt{add-flight-import.tsx}) rufen \texttt{addFlight} bzw.\ \texttt{updateFlight} direkt am Store auf.
Der Store berechnet fehlende Werte (Distanz, Dauer) und delegiert anschließend an
den \texttt{supabaseService}:

\begin{lstlisting}[caption={Datenfluss Flug anlegen (vereinfacht)},label={lst:flight_dataflow},language={}]
UI-Screen (ensureAirportByCode vor Aufruf)
  -> store.addFlight(flightData)
        | Distanz berechnen (haversineDistanceKm)
        | Dauer berechnen (computeDurationMinutesFromTimestamps)
        -> supabaseService.createFlight(...)
              | Airport via getOrCreateAirport() in DB persistieren
              | INSERT in user_flights (mit RLS)
              -> Rueckgabe: vollstaendiges Flight-Objekt (inkl. Airport-Join)
\end{lstlisting}

\subsection{Distanz- und Dauerberechnung}

Distanz und Flugdauer werden im Store automatisch ermittelt, sofern sie nicht
bereits durch den Import geliefert wurden. Für die Distanz kommt die
Haversine-Formel zum Einsatz (vgl.\ auch Kapitel~\ref{chapter:weltkarte}); die Dauer ergibt sich aus der
Differenz von \texttt{departureAt} und \texttt{arrivalAt}:

\begin{lstlisting}[style=skyline-ts,caption={Distanz- und Dauerberechnung im Store (Auszug store/index.ts)},label={lst:addflight_distance_duration}]
if (!distanceStr && flightData.from?.latitude != null
                 && flightData.to?.latitude != null) {
  distanceKm = haversineDistanceKm(flightData.from, flightData.to);
  distanceStr = formatDistanceKm(distanceKm);
}

if (!durationStr && flightData.departureAt && flightData.arrivalAt) {
  const mins = computeDurationMinutesFromTimestamps(
    flightData.departureAt, flightData.arrivalAt
  );
  const h = Math.floor(mins / 60);
  const m = mins % 60;
  durationStr = h > 0 && m > 0 ? `${h}h ${m}m`
              : h > 0            ? `${h}h`
              :                    `${m}m`;
}
\end{lstlisting}

\subsection{Status-Berechnung}

Der Flugstatus (\texttt{upcoming} / \texttt{completed}) wird beim Speichern automatisch aus
\texttt{departureAt} relativ zur aktuellen Zeit berechnet und mit \texttt{getEffectiveFlightStatus}
beim Update normalisiert -- ohne manuelle Eingabe durch den Nutzer.

\subsection{Airport-Auflösung}

Airports werden über die Funktion \texttt{ensureAirportByCode()} aufgelöst. Dabei wird
zunächst über den \texttt{AirportAPI}-Service versucht, den Airport per IATA- oder
ICAO-Code zu finden. Liefert die API (einschließlich der vordefinierten
\texttt{POPULAR\_AIRPORTS}) keinen Treffer, wird \emph{kein} neuer Airport in der
lokalen \texttt{airports}-Tabelle angelegt und die Funktion gibt \texttt{null} zurück.

Nur wenn ein Airport von der API bzw.\ aus \texttt{POPULAR\_AIRPORTS} gefunden wird,
wird er über \texttt{upsertAirportFromApi()} in der Datenbank angelegt oder -- falls
bereits vorhanden -- wiederverwendet. Dadurch bleibt das Datenmodell konsistent,
gleichzeitig wird aber verhindert, dass für komplett unbekannte oder fehlerhafte
Codes fiktive Airports angelegt werden.

% ============================================================
\section{Dokumente, Notizen und Checklisten}
% ============================================================

Reisende können Flugdokumente (Boarding-Pass, Buchungsbestätigung, Quittungen) über
den \texttt{documentService} hochladen. Dateiinhalt landet im Storage-Bucket
\texttt{flight-documents}; Metadaten (Dateiname, MIME-Type, Größe, Storage-Pfad,
Signed URL) werden in \texttt{flight\_documents} referenziert. Der Zugriff erfolgt
ausschließlich über zeitlich begrenzte Signed URLs, die bei Ablauf automatisch
erneuert werden. Details zu Upload-Ablauf und Rollback-Logik sind in Kapitel~\ref{chapter:datenverwaltung}
dokumentiert.

Freitextliche Notizen (\texttt{user\_notes}) und Checklisten (\texttt{user\_checklists},
\texttt{user\_checklist\_items}) sind über \texttt{flight\_id} an den jeweiligen Flug gebunden und
durch RLS-Policies auf den jeweiligen Nutzer beschränkt. Checklisten können optional
mit einem Erinnerungszeitpunkt (\texttt{reminder\_at}) versehen werden; der
\texttt{noteChecklistReminderService} plant in diesem Fall eine lokale Notification (vgl.\
Kapitel~\ref{chapter:notifications}). Die Trip-Detailansicht lädt Notizen und Checklisten beim Öffnen parallel
via \texttt{Promise.allSettled}, sodass keine sequenziellen Wartezeiten entstehen.

% ============================================================
\section{Statistiken}
% ============================================================

Statistiken werden clientseitig aus den im Store vorliegenden Flugdaten berechnet --
es gibt keine separaten Aggregationsabfragen an die Datenbank. Die Berechnung
erfolgt in einem \texttt{useMemo}-Hook, der nur neu berechnet wird, wenn sich die
Flugdaten ändern.

\subsection{Nutzer-Statistiken}

Für das Nutzerprofil werden folgende Kennzahlen berechnet:

\begin{table}[H]
  \centering
  \caption{Nutzer-Statistiken (Kennzahlen)}
  \label{tab:user_stats}
  \begin{tabular}{p{4cm}p{9cm}}
  \toprule
  \textbf{Kennzahl} & \textbf{Beschreibung} \\
  \midrule
  \texttt{tripsTaken} & Anzahl aller Flüge des Nutzers \\
  \texttt{totalDistance} & Summierte Flugdistanz in km \\
  \texttt{countriesVisited} & Eindeutige Länder aus \texttt{from.country} und \texttt{to.country} \\
  \texttt{favoriteDestination} & Häufigstes Reiseziel (nach \texttt{to.city}) \\
  \texttt{averageTripDuration} & Durchschnittliche Flugdauer in Stunden \\
  \texttt{nextFlight} & Nächster bevorstehender Flug (sortiert nach \texttt{departureAt}) \\
  \texttt{currentStreak} & Aktueller Reise-Streak (max.\ 5, gamifiziert) \\
  \bottomrule
  \end{tabular}
\end{table}

\begin{lstlisting}[style=skyline-ts,caption={Berechnung totalDistance und countriesVisited (Auszug profile.tsx)},label={lst:user_stats_code}]
const totalDistance = userFlights.reduce((sum, f) => {
  if (!f.distance) return sum;
  const match = f.distance.replace(/[^\d]/g, '').match(/(\d+)/);
  const num = match ? parseInt(match[1], 10) : 0;
  return sum + (isNaN(num) ? 0 : num);
}, 0);

const countriesVisited = new Set<string>();
userFlights.forEach(f => {
  if (f.from?.country) countriesVisited.add(f.from.country);
  if (f.to?.country)   countriesVisited.add(f.to.country);
});
\end{lstlisting}

\subsection{Company-Statistiken}

Für die aktive Firma (sichtbar für Owner) werden Gesamtflüge, Gesamtdistanz und
Länderanzahl aller Firmenmitglieder aggregiert. Die Abfrage läuft direkt gegen
\texttt{user\_flights} gefiltert auf die \texttt{profile\_id}-Liste aller \texttt{company\_members} und die
\texttt{company\_id}:

\begin{lstlisting}[style=skyline-ts,caption={Company-Statistiken (Auszug app/company/index.tsx)},label={lst:company_stats}]
const { data } = await supabase
  .from('user_flights')
  .select(`id, distance_km, distance,
           from_airport:from_airport_id (country),
           to_airport:to_airport_id (country)`)
  .in('profile_id', memberIds)
  .eq('company_id', activeMembership.companyId);
\end{lstlisting}

Ein CSV-Export für Flugdaten und Company-Statistiken ist als Erweiterung vorgesehen und soll die Weitergabe an externe Abrechnungssysteme ermöglichen.

% ============================================================
\section{Company-Features}
% ============================================================

\subsection{Firmenstruktur und Rollen}

Firmen (\texttt{companies}) werden von einem Owner angelegt. Mitglieder können die Rolle
\texttt{owner} oder \texttt{worker} innehaben; die Mitgliedschaft wird in \texttt{company\_members}
verwaltet. RLS-Policies erlauben Lesezugriff auf Profile innerhalb derselben Firma.
Gemeinsame Firmenflüge werden über \texttt{user\_flights.company\_id} und die Mitgliedschaft in \texttt{company\_members} abgebildet; Flüge mit gesetzter \texttt{company\_id} gehören zur Firma, der Zugriff erfolgt über die Mitglieder.

\subsection{Invite-Code-System}

Das Beitreten zu einer Firma erfolgt über einen 8-stelligen Einladungscode. Der Code
wird aus dem Alphabet \texttt{ABCDEFGHJKLMNPQRSTUVWXYZ23456789} (ohne leicht
verwechselbare Zeichen wie \texttt{0}, \texttt{O}, \texttt{1}, \texttt{I}) zufällig generiert und auf
Einzigartigkeit in der \texttt{companies}-Tabelle geprüft (max.\ 10 Versuche). Der Code hat
kein Ablaufdatum und bleibt dauerhaft gültig.

\begin{lstlisting}[style=skyline-ts,caption={Generierung eindeutigen Invite-Codes (Auszug companyService.ts)},label={lst:invite_code}]
const INVITE_CODE_ALPHABET = 'ABCDEFGHJKLMNPQRSTUVWXYZ23456789';
const INVITE_CODE_LENGTH = 8;

function generateInviteCode(): string {
  let code = '';
  for (let i = 0; i < INVITE_CODE_LENGTH; i += 1) {
    code += INVITE_CODE_ALPHABET.charAt(
      Math.floor(Math.random() * INVITE_CODE_ALPHABET.length)
    );
  }
  return code;
}

async function generateUniqueCompanyInviteCode(): Promise<string> {
  for (let i = 0; i < 10; i++) {
    const candidate = generateInviteCode();
    const { data } = await supabase
      .from('companies')
      .select('id')
      .eq('invite_code', candidate)
      .maybeSingle();
    if (!data) return candidate;
  }
  throw new Error('Unable to generate invite code');
}
\end{lstlisting}

\subsection{Beitrittsprozess}

Beim Einlösen eines Codes wird dieser normalisiert (trim + toUpperCase), in
\texttt{companies} gesucht und -- sofern der Nutzer noch kein Mitglied ist -- ein Eintrag
mit \texttt{role: 'worker'} in \texttt{company\_members} eingefügt. Ist der Nutzer bereits
Mitglied, wird die bestehende Rolle zurückgegeben ohne einen Fehler zu werfen.

\begin{lstlisting}[style=skyline-ts,caption={Beitritt per Invite-Code (Auszug companyService.ts)},label={lst:join_invite}]
const code = codeInput.trim().toUpperCase();
const { data: company } = await supabase
  .from('companies')
  .select('id')
  .eq('invite_code', code)
  .maybeSingle();

if (!company) throw new Error('Invite code not found');

const { data: membership } = await supabase
  .from('company_members')
  .select('role')
  .eq('company_id', company.id)
  .eq('user_id', currentUser.id)
  .maybeSingle();

if (membership) return { companyId: company.id, joinedAs: membership.role };

await supabase
  .from('company_members')
  .insert({ company_id: company.id, user_id: currentUser.id, role: 'worker' });
\end{lstlisting}

% ============================================================
\section{Gamification}
% ============================================================

Achievements motivieren Nutzer durch definierte Meilensteine (z.\,B.\ erste Flugreise,
bestimmte zurückgelegte Gesamtstrecke). Die globalen Achievement-Definitionen liegen
in der Tabelle \texttt{achievements} und sind öffentlich lesbar, da sie keine
personenbezogenen Daten enthalten. Freigeschaltete Achievements werden pro Nutzer in
\texttt{user\_achievements} mit \texttt{unlocked\_at} und \texttt{metadata} gespeichert.

Das System ist bewusst einfach gehalten: Es gibt keine XP-Punkte, kein Levelsystem
und keine Rangliste. Beim Freischalten eines Achievements erhält der Nutzer eine
Toast-Nachricht und haptisches Feedback als unmittelbare Rückmeldung.

% ============================================================
\section{Teststrategie \& Validierung}
% ============================================================

Skyline wurde primär auf iOS-Geräten getestet. Manuelle UI-Tests sichern die
Hauptabläufe (Flug hinzufügen, Import, Trip-Details, Dokumente, Notizen,
Checklisten). Reminder-Tests prüfen Offsets, Quiet Hours und Reschedule-Verhalten
nach App-Neustart; Import-Tests decken QR-Scan (BCBP), OCR und
E-Mail-Screenshot-Upload ab. Statistiken und Achievements werden bei typischen
Nutzerabläufen mitgeprüft. Die Android-Variante wurde im Projektzeitraum nur
eingeschränkt validiert; eine vollständige Android-Testabdeckung ist als
Erweiterung vorgesehen.

\begin{figure}[H]
  \centering
  \fbox{\parbox{0.85\textwidth}{\centering\small Screenshot Testfälle / UI-Flows -- hier einfügen.}}
  \caption{Testfälle und UI-Flows (Platzhalter)}
  \label{fig:test_ui_flows}
\end{figure}

% ============================================================
\section{Deployment \& App Store Release}\label{sec:deployment-appstore}
% ============================================================

Dieser Abschnitt dokumentiert den Weg von Skyline vom Build bis zur Veröffentlichung
im Apple App Store: verwendete Build-Konfiguration, Einreichung über App Store Connect,
Nutzung von TestFlight sowie Erfahrungen im Review-Prozess.

\subsection{Build-Konfiguration (EAS Build)}

Die iOS-Builds werden mit \textbf{EAS Build} (Expo Application Services) erzeugt
\cite{easBuild}. Konkret wurde ein \textbf{Produktions-Profil} (\emph{production})
für die App-Store-Version verwendet. Ein separates Build-Schema wurde nicht
definiert; die Konfiguration erfolgt über \texttt{eas.json}. Sensible Werte wie
API-Keys werden über Environment-Variablen in der lokalen \texttt{.env}-Datei
verwaltet.

\subsection{App Store Connect \& Einreichung}

Über \textbf{App Store Connect} wurde die App bei Apple eingereicht. Die konkreten
Angaben: \textbf{Bundle ID} \texttt{com.borisplesnicar.skyline}, App-Name
\textbf{Skyline Hub}, \textbf{Version} 2.2.0 (Build 14), \textbf{Kategorie} Travel.
Die Store-Beschreibung lautet: \emph{Skyline is an app that lets you track and
organize your flights and travels. You can add flights, view your past journeys
on an interactive world map, and see simple statistics such as countries visited
or total distance flown. In addition, you can store important travel documents
and photos for each trip and monitor your progress through goals and achievements.}
Datenschutzangaben sind in der App hinterlegt; Support-URL und Nutzungsbedingungen
verweisen auf \url{https://skylinehub.vercel.app/}. Für den Store wurden
Screenshots im Format iPhone 6,5'' Display eingereicht (Abbildung~\ref{fig:store_screenshots}).

\begin{figure}[H]
  \centering
  \includegraphics[width=0.48\textwidth]{screens/Onboarding_01}\hfill
  \includegraphics[width=0.48\textwidth]{screens/Onboarding_02}
  \caption{Store-Screenshots (iPhone 6,5'') -- Auswahl der eingereichten App-Store-Bilder}
  \label{fig:store_screenshots}
\end{figure}

\subsection{TestFlight für interne Tests}

\textbf{TestFlight} ermöglicht die Verteilung von Builds an interne und externe
Tester vor der öffentlichen Freigabe \cite{testflight}. TestFlight wurde für
interne Tests genutzt: neben den Entwicklern wurden auch Klassenkollegen als
Tester eingeladen, um die App vor der Veröffentlichung zu validieren
\cite{ralfebertTestflight}.

\subsection{Review-Prozess und Erfahrungen}

Apple prüft jede App-Einreichung nach den App-Store-Richtlinien; der Review kann
mehrere Tage dauern und bei Beanstandungen zu Ablehnung oder Nachfragen führen
\cite{appleAppReview}. Die Einreichung von Skyline Hub wurde zunächst
\emph{abgelehnt}: Apple monierte fehlende Mehrsprachigkeit sowie das Fehlen eines
QR-Codes, mit dem die Scan-Funktion (BCBP-Barcode) von den Prüfern getestet werden
konnte. Nach Behebung dieser Punkte wurde die App in einem weiteren Review-Durchlauf
angenommen. Abbildung~\ref{fig:review_verlauf} zeigt den Verlauf der Einreichung
bis zur Freigabe.

\begin{figure}[H]
  \centering
  \includegraphics[width=0.9\textwidth]{verlauf-rueckfragen}
  \caption{Verlauf der App-Store-Einreichung bis zur Annahme}
  \label{fig:review_verlauf}
\end{figure}

\subsection{App im Store}

Skyline Hub ist im Apple App Store verfügbar:
\url{https://apps.apple.com/us/app/skyline-hub/id6758091952}

% ============================================================
\section{Überblick über die App und Benutzeroberfläche}\label{sec:app-overview}
% ============================================================

Dieser Abschnitt gibt einen Überblick über die Benutzeroberfläche von Skyline.
Features, die in den Implementierungskapiteln~5--8 ausführlich behandelt werden
(Import, Weltkarte, Benachrichtigungen, Datenverwaltung), werden hier nur kurz
referenziert. Für alle übrigen Features -- Home-Screen, Navigation, Statistik-Ansicht,
Gamification und Company-Features -- ist dies die einzige Stelle in der Arbeit, an der
sie mit Screenshots und Sinn-Erklärung dokumentiert werden.

\subsection{Forschungsmodule (Kurzreferenz)}

Die vier Forschungsmodule werden in eigenen Kapiteln detailliert beschrieben:

\begin{itemize}
  \item \textbf{Import} (Kapitel~\ref{chapter:import}): OCR-basierte Erkennung von
    Buchungsbestätigungen aus Screenshots und PDFs, BCBP-Barcode-Scan, manuelle
    Flugeingabe mit Airport-Autocomplete.
  \item \textbf{Weltkarte} (Kapitel~\ref{chapter:weltkarte}): Interaktive Kartenansicht
    mit Great-Circle-Routen, Flugvisualisierung und Distanzberechnung.
  \item \textbf{Benachrichtigungen} (Kapitel~\ref{chapter:notifications}): Lokale
    Erinnerungen für Checklisten und Notizen, Quiet Hours, Reschedule nach App-Neustart.
  \item \textbf{Datenverwaltung} (Kapitel~\ref{chapter:datenverwaltung}): RLS-Architektur,
    Dokumenten-Upload, Notizen und Checklisten pro Flug.
\end{itemize}

\subsection{Home-Screen und Navigation}

Der Home-Screen ist die zentrale Einstiegsseite nach dem Login. Er zeigt eine
Übersicht der nächsten Flüge, eine Schnellaktion zum Hinzufügen neuer Flüge sowie
-- bei Firmenmitgliedschaft -- einen direkten Zugang zum Company-Workspace. Die
Gesamtdistanz und der aktuelle Reise-Streak werden prominent angezeigt und
motivieren zur regelmäßigen Nutzung.

Die Navigation erfolgt über eine \textbf{floating Tab-Bar} mit vier Tabs: Home,
Karte, Profil und Einstellungen. Die Tab-Bar ist als schwebendes, abgerundetes
Element am unteren Bildschirmrand positioniert und nutzt einen animierten
Sliding-Indicator zur Hervorhebung des aktiven Tabs.

\begin{figure}[H]
  \centering
  \fbox{\parbox{0.75\textwidth}{\centering\small
    \textbf{[Screenshot: Home-Screen mit Flugübersicht und Tab-Bar einfügen -- z.\,B.\ \texttt{content/img/home\_screen.png}]}
  }}
  \caption{Home-Screen von Skyline mit Flugübersicht und floating Tab-Bar}
  \label{fig:home_screen}
\end{figure}

\subsection{Statistik-Ansicht und Profil}

Das Profil zeigt die persönlichen Reise-Statistiken:
Anzahl der Flüge, zurückgelegte Gesamtstrecke, besuchte Länder, beliebtestes Reiseziel,
durchschnittliche Flugdauer und den nächsten bevorstehenden Flug. Die Berechnung
erfolgt clientseitig aus den im Store vorliegenden Flugdaten (vgl.\ Abschnitt~4.5).
Zusätzlich wird der aktuelle \emph{Reise-Streak} angezeigt -- ein Gamification-Element,
das Vielreisende dazu anregen soll, die App kontinuierlich zu nutzen und Flüge zeitnah
einzutragen. Über das Profil gelangt der Nutzer auch zu Achievements, Einstellungen
und dem Company-Workspace.

\begin{figure}[H]
  \centering
  \fbox{\parbox{0.75\textwidth}{\centering\small
    \textbf{[Screenshot: Profil mit Statistik-Kennzahlen einfügen -- z.\,B.\ \texttt{content/img/profile\_stats.png}]}
  }}
  \caption{Profil-Screen mit Nutzer-Statistiken und Streak-Anzeige}
  \label{fig:profile_stats}
\end{figure}

\subsection{Gamification und Achievements}

Achievements werden bei definierten Meilensteinen freigeschaltet (z.\,B.\ erster Flug,
10 Flüge, 5 Länder, 10.000\,km Gesamtstrecke). Der \textbf{Sinn} dahinter: Vielreisende
sollen motiviert werden, die App aktiv zu nutzen und Reisedaten vollständig
einzupflegen. Vollständige Daten verbessern wiederum die Qualität der Statistiken
und -- bei Firmenmitgliedschaft -- die Aussagekraft der Company-Übersicht. Das System
ist bewusst schlank gehalten: keine XP-Punkte, kein Levelsystem, keine Rangliste.
Beim Freischalten erscheint eine Toast-Nachricht und haptisches Feedback.

Die Achievements-Ansicht listet alle verfügbaren Meilensteine mit Fortschrittsbalken.
Bereits freigeschaltete Achievements werden visuell hervorgehoben.

\begin{figure}[H]
  \centering
  \fbox{\parbox{0.75\textwidth}{\centering\small
    \textbf{[Screenshot: Achievements-Übersicht mit Fortschrittsbalken einfügen -- z.\,B.\ \texttt{content/img/achievements.png}]}
  }}
  \caption{Achievements-Screen mit Meilensteinen und Fortschrittsanzeige}
  \label{fig:achievements}
\end{figure}

\subsection{Company-Features und Firmenansicht}

Company-Features richten sich an Firmen, die Dienstreisen ihrer Mitarbeiter zentral
erfassen und auswerten möchten. Der \textbf{Business-Sinn}: Die Firma erhält einen
Überblick über alle Dienstreisen der Mitarbeiter, vereinfacht die Abrechnung und
kann Reisekosten und -muster analysieren. Flüge werden beim Anlegen optional einer
Firma zugeordnet (\texttt{company\_id}); nur Mitglieder derselben Firma sehen die
aggregierten Daten.

Die Firmenansicht zeigt für Owner: Gesamtflüge,
Gesamtdistanz und Länderanzahl aller Mitglieder, die Einladung von neuen Mitgliedern
über einen Invite-Code sowie die Liste der Firmenflüge. Worker sehen die aggregierten
Statistiken und können bei mehreren Firmenmitgliedschaften zwischen diesen wechseln.
Der Beitritt erfolgt über einen 8-stelligen Einladungscode (vgl.\ Abschnitt~4.6).

\begin{figure}[H]
  \centering
  \fbox{\parbox{0.75\textwidth}{\centering\small
    \textbf{[Screenshot: Company-Übersicht mit Statistiken und Mitgliederliste einfügen -- z.\,B.\ \texttt{content/img/company\_overview.png}]}
  }}
  \caption{Company-Screen mit aggregierten Firmen-Statistiken}
  \label{fig:company_overview}
\end{figure}

\subsection{Trip-Details und Flug hinzufügen}

Die Trip-Detailansicht zeigt einen einzelnen Flug mit allen zugehörigen Daten:
Abflug/Ziel, Zeiten, Dokumente, Notizen und Checklisten. Das Hinzufügen neuer
Flüge erfolgt über manuelle Eingabe oder Import aus Screenshot/PDF/BCBP; Details
zum Import-Flow sind in Kapitel~\ref{chapter:import} dokumentiert.

\subsection{Flug-Kalender}

Der Flug-Kalender bietet eine monatsbasierte Übersicht aller Flüge und
Erinnerungstermine. Nutzer können zwischen Flügen und geplanten Events
(Checklisten-Reminder, Notizen-Erinnerungen) wechseln oder beide Ansichten
kombiniert anzeigen. Ein Tipp auf ein Datum zeigt die Flüge und Termine des
Tages; von dort aus gelangt man direkt zur Trip-Detailansicht. Der Kalender
unterstützt die Planung von Reisen im Zeitverlauf und macht anstehende
Erinnerungen auf einen Blick sichtbar.

\begin{figure}[H]
  \centering
  \fbox{\parbox{0.75\textwidth}{\centering\small
    \textbf{[Screenshot: Flug-Kalender mit Monatsansicht einfügen]}
  }}
  \caption{Flug-Kalender mit Monatsübersicht und Tagesdetails}
  \label{fig:flight_calendar}
\end{figure}
